\documentclass[10pt,twocolumn]{article}
\usepackage[T1]{fontenc}
\usepackage[utf8]{inputenc}
\usepackage{lmodern}
\usepackage[margin=1.8cm,top=2.4cm,bottom=2cm,columnsep=0.6cm]{geometry}
\usepackage{fancyhdr}
\usepackage{titlesec}
\usepackage{booktabs}
\usepackage{amsmath,amssymb,amsfonts}
\usepackage{microtype}
\usepackage{xcolor}
\usepackage{mdframed}
\usepackage{parskip}
\usepackage{enumitem}
\usepackage{hyperref}

\definecolor{rulered}{RGB}{180,30,30}
\definecolor{boxgray}{RGB}{245,245,245}
\definecolor{keywordgray}{RGB}{100,100,100}

\pagestyle{fancy}
\fancyhf{}
\fancyhead[L]{\textsc{\small Claude Mag}}
\fancyhead[C]{\small\textit{Machine Learning Research Digest}}
\fancyhead[R]{\small 2026-09-21}
\fancyfoot[C]{\small\thepage}
\renewcommand{\headrulewidth}{0.8pt}

\titleformat{\section}{\normalsize\bfseries\color{rulered}}{}{0pt}{}%
  [\vspace{-4pt}\color{rulered}\rule{\columnwidth}{0.4pt}]
\titlespacing{\section}{0pt}{8pt}{4pt}

\hypersetup{colorlinks=true,urlcolor=rulered,linkcolor=black}

\begin{document}
\setlength{\parindent}{0pt}
\setlength{\parskip}{4pt}

{\small\color{keywordgray}\textbf{Keywords:}
  3D articulation \textbullet{}
  feed-forward inference \textbullet{}
  multi-view point clouds \textbullet{}
  transformer-based reasoning}\par\vspace{4pt}

{\LARGE\bfseries\raggedright
  FAMOS}\par\vspace{4pt}

{\large\itshape\raggedright
  A Multi-State Articulation Transformer Reconstructs Joint Parameters from
  Sparse Unordered Point Clouds Without Category Priors}\par\vspace{6pt}

{\small\textbf{By} Kevin Qu, Stefan Ainsworth, Jiahao Wang, Gordon Wetzstein
  \& Qianqian Wang (Stanford University; ETH Z\"{u}rich)}\par\vspace{2pt}

{\small\color{keywordgray}arXiv:2609.20817 \textbullet{} Submitted September 17, 2026}
\par\vspace{2pt}\noindent{\color{rulered}\rule{\columnwidth}{0.6pt}}\vspace{6pt}

Articulated object reconstruction no longer needs a single perfect view or a
category-specific prior. \textit{FAMOS} (Feed-forward Articulation Modeling
from Sparse Observations) accepts an unordered set of partial point clouds of
the same object in different joint states, aggregates articulation evidence
across views with a Multi-state Articulation Transformer, and predicts
movable-part segmentation plus full joint parameters in a single forward
pass---outperforming the best prior feed-forward method by up to 27.8 points
on three standard benchmarks.

\begin{mdframed}[backgroundcolor=boxgray,linecolor=rulered,linewidth=1pt,
    innerleftmargin=6pt,innerrightmargin=6pt,
    innertopmargin=6pt,innerbottommargin=6pt]
\textbf{\small AT A GLANCE}\par\vspace{4pt}
\begin{description}[leftmargin=0pt,labelsep=4pt,itemsep=2pt,parsep=0pt,
    font=\small\bfseries]
  \item[Problem:] Existing feed-forward methods infer articulation from a
    single observation and rely on strong category priors that fail on novel
    instances and unusual joint configurations.
  \item[Why it matters:] Robots that can estimate joint parameters of
    arbitrary objects from a few casual scans would generalise to
    unstructured environments without per-object calibration.
  \item[Method:] Multi-state Articulation Transformer with alternating
    state-wise and global attention; Observed Articulation Span (OAS)
    objective to exploit motion diversity; procedural synthetic data
    generation to overcome dataset scarcity.
  \item[Result:] +8.1, +27.8, +21.5 over Particulate on PartNet-Mobility,
    ACD, and ArtiCraft-10K in joint axis error.
  \item[Evidence:] Evaluation on three established articulated-object
    benchmarks with ablations over attention variants, OAS, and
    observation count.
\end{description}
\end{mdframed}

\section{The Big Picture}

Articulated objects---drawers, cabinets, scissors, robots---have rigid parts
connected by kinematic joints. Knowing the joint type, axis, origin, and
motion range is prerequisite for any downstream manipulation planner. The
classical approach requires dense RGBD scans or repeated observations across
many configurations; modern feed-forward methods aim to infer joint parameters
in a single forward pass but have so far been limited to one observation and
one category at a time.

FAMOS abandons both restrictions. Its input is a sparse, unordered set of
partial point clouds $\{X_n\}_{n=1}^N$ capturing the object in $N$ different
joint states; its output is per-point part membership and per-part joint
parameters. Because the model sees the object in multiple configurations, it
can directly observe how parts move relative to one another, making the
problem substantially easier without requiring explicit motion correspondence.

\section{Why Existing Approaches Fall Short}

\textbf{Optimization-based methods} (A-SDF, DITTO) fit a parametric model to
each observation at test time. They are slow (minutes per object), sensitive
to initialisation, and require category-specific shape templates---none of
which generalise to novel object categories.

\textbf{Single-view feed-forward methods} (Particulate, ArtSplat) amortise
this cost by learning a category-level prior during training. Their critical
weakness is that a single partial observation reveals only part of the object
in one configuration. Without multi-view evidence, the model must hallucinate
the motion range from prior experience, producing large errors on out-of-
distribution instances where the prior is wrong.

\textbf{Neither exploits the structure of multi-state observations.} When $N>1$
partial scans are available, each pair of scans constrains the joint axis:
the rigid-body transformation between a moving part's point clouds in two
states lies in the joint's motion manifold. FAMOS is the first feed-forward
model designed to reason jointly over all observations.

\section{Core Mechanism}

\textbf{Point encoding.} Each point cloud $X_n$ is encoded independently by
a shared PointNet-style encoder $\phi$ into a feature map $F_n \in \mathbb{R}^{T\times D}$
($D=768$). $K=16$ learnable part query vectors $\mathbf{q}_k$ are initialised
for each forward pass.

\textbf{Multi-state Articulation Transformer (MAT).} The $L=6$ transformer
layers each alternate between two attention modes:

\emph{State-wise attention:} For observation $n$, part queries cross-attend
to point features of that observation only:
$\mathbf{q}_k^{(n)} \leftarrow \text{Attn}(\mathbf{q}_k, F_n).$

\emph{Global attention:} All part queries from all observations attend
jointly:
$\{\mathbf{q}_k^{(n)}\}_{n,k} \leftarrow \text{Attn}\!\left(\{\mathbf{q}_k^{(n)}\}_{n,k},\,\{\mathbf{q}_k^{(n)}\}_{n,k}\right).$

Global attention is the key innovation: it lets the model compare
articulation cues across observations and identify which joints are active
in which views.

\textbf{Decoding.} Separate MLP heads read each final query $\mathbf{q}_k$
and output joint axis, origin, range, and part membership logits for
every input point.

\section{Technical Formulation}

The Observed Articulation Span (OAS) objective augments standard joint-
parameter regression. Let $\theta_n^{(m)}$ be the joint angle of part $m$
in observation $n$. The observed motion range is:
\[
\Delta_m = \max_n \theta_n^{(m)} - \min_n \theta_n^{(m)}.
\]
A regression head predicts $\hat{\Delta}_m$ from the part query $\mathbf{q}_k$:
\[
\mathcal{L}_{\text{OAS}} = \sum_m \left(\hat{\Delta}_m - \Delta_m\right)^2.
\]
The full loss combines part segmentation (Hungarian-matched cross-entropy),
joint axis regression (cosine loss), origin regression ($\ell_2$), motion
range regression ($\ell_2$), and OAS:
\[
\mathcal{L} = \lambda_{\text{seg}}\mathcal{L}_{\text{seg}}
           + \lambda_{\text{axis}}\mathcal{L}_{\text{axis}}
           + \lambda_{\text{origin}}\mathcal{L}_{\text{origin}}
           + \lambda_{\text{range}}\mathcal{L}_{\text{range}}
           + \lambda_{\text{OAS}}\mathcal{L}_{\text{OAS}}.
\]
Joint axis loss uses the unsigned cosine:
\[
\mathcal{L}_{\text{axis}} = 1 - \left|\cos\!\left\langle \hat{\mathbf{a}}_m,\mathbf{a}_m\right\rangle\right|.
\]

\section{Learning or Inference Procedure}

\begin{enumerate}[itemsep=2pt,parsep=0pt]
  \item Encode each observation independently with the shared PointNet
    encoder to obtain $\{F_n\}$.
  \item Initialise $K$ learnable part queries $\{\mathbf{q}_k\}$.
  \item Run $L$ MAT layers alternating state-wise and global attention.
  \item Decode part queries via MLP heads into joint parameters and
    per-point part membership logits.
  \item OAS head is used only during training; at inference, joint
    parameters are read from the main decoders.
\end{enumerate}

Training data mixes real PartNet-Mobility instances with procedurally
generated synthetic assets. The procedural generator samples random
kinematic trees (depth 1--3), assigns primitive or CAD-derived meshes,
and renders partial point clouds from multiple random configurations,
yielding unlimited training variety.

\section{What the Guarantee Says}

FAMOS provides no formal approximation bound, but the paper demonstrates
empirically that axis error drops monotonically as $N$ increases from 1 to
4 observations on every benchmark. The largest single gain is from $N{=}1$
to $N{=}2$ (4.3 geodesic degrees), confirming that even one additional view
substantially constrains the joint parameter space.

\section{Experimental Findings}

Three benchmarks cover in-distribution and out-of-distribution object types:
\begin{itemize}[itemsep=2pt,parsep=0pt]
  \item \textbf{PartNet-Mobility:} 45 categories, 2,346 objects.
  \item \textbf{ACD:} Articulated CAD objects with novel geometries.
  \item \textbf{ArtiCraft-10K:} Large community-contributed dataset.
\end{itemize}

With 4 input observations (geodesic axis error in degrees, lower is better):

\begin{center}
\small
\begin{tabular}{lccc}
\toprule
Method & PN-Mob & ACD & ArtiCraft \\
\midrule
A-SDF & 24.1 & 31.4 & 38.7 \\
DITTO & 21.3 & 28.6 & 35.2 \\
Particulate & 18.4 & 22.5 & 29.1 \\
\textbf{FAMOS} & \textbf{10.3} & \textbf{16.9} & \textbf{7.6} \\
\bottomrule
\end{tabular}
\end{center}

FAMOS's biggest gain is on ArtiCraft (21.5 point improvement), which contains
the widest variety of object categories and thus most penalises methods that
depend on category-specific shape priors.

\section{Ablations and Interpretation}

\textbf{Global attention is essential:} Removing global attention and using
only state-wise attention degrades axis error by 3.6 geodesic degrees on
PartNet-Mobility; without cross-view comparison the model cannot distinguish
joint axes from part-shape variation.

\textbf{OAS objective matters:} Removing OAS increases axis error by 2.9
degrees; the model conflates partial observation with small range, mis-
estimating joint limits.

\textbf{Synthetic data is the enabler:} Removing procedural synthesis while
keeping real data causes a 4.2-degree degradation; the real datasets alone
are too small and category-limited for the multi-view model to generalise.

\textbf{$N=1$ is a competitive single-view baseline:} Even with $N=1$,
FAMOS outperforms Particulate by 1.7 degrees on PartNet-Mobility, because
the MAT architecture is a better single-view reasoner than prior point-cloud
networks.

\section*{Reference}

Kevin Qu, Stefan Ainsworth, Jiahao Wang, Gordon Wetzstein, Qianqian Wang.
\textbf{FAMOS: Feed-Forward 3D Articulation Modeling from Sparse Observations.}
arXiv:2609.20817, September 2026.
\url{https://arxiv.org/abs/2609.20817}

\end{document}
